\section{Analysis and Discussion}
\label{sec:discussion}

\subsection{System-Level Performance}

The seven-benchmark suite shows one runtime operating across software repair,
GPU-kernel optimization, model training, research tasks, and data synthesis; a
separate mathematical vertical trace examines proof-oriented research depth. On
SWE-Bench Pro, \systemname reaches approximately 78\%
accuracy compared with 59\% for Direct Copilot at 1.41$\times$ aggregate Tokens.
The comparison pairs additional computation for planning, execution, and review
with a 19-point accuracy difference.

The longitudinal run shows how this cost changes with accumulated state. Mature
W19--22 uses 21\% fewer solve input Tokens and 15\% less active time per task than
startup W1--6. The curve remains non-monotone: W13--18 combines low Token use with
longer execution, while W23--24 contains a late concentration of difficult tasks. The operating profile changes
over time as shared state expands, while task-dependent variation remains visible.

\subsection{Reviewer as Selective Error Correction}

The Reviewer requests revision on 43 SWE-Bench Pro tasks. After another Engineer
round, 34 pass the official verifier and 22 complete the stricter
\code{continue}$\rightarrow$revision$\rightarrow$\code{done} loop. These trajectories
show the Reviewer functioning as an error-correction stage rather than a terminal
commentary layer.

Reviewer routing is adaptive: Reviewer-routed tasks use more Tokens and active time
than self-reviewed tasks. The observed routing concentrates independent review on
the more resource-intensive part of the workload.

\subsection{Autonomous Research Production}

The six-paper case study extends the endpoint results into a complete research
lifecycle. All six canonical pipelines reach submission completion across six
domains and two venue formats, but none follows a one-pass path. Reviewer revision
verdicts, Stage rollbacks, and session rolls are common rather than exceptional.
The final manuscripts therefore demonstrate persistence across research framing,
experimentation, writing, review, and packaging---not merely the ability to emit
paper-like prose.

The multimodal-hallucination trace illustrates the control mechanism. Reviewer
gates reject seven weak method routes before they are scaled into an unsupported
positive claim; Planner proposes reframing the accumulated failures as a diagnostic study;
Engineer then completes the 15-cell canonical matrix; and Manager accepts two late
rollbacks when submission checks expose scorer-provenance and telemetry defects.
The recorded sequence is consistent with role separation redirecting the scientific
trajectory: failure is
retained as information, while only reviewed state controls the next Stage.

The case also separates artifact completion from process certification. Five final
assurance snapshots pass; the compositional-matching project retains a stale blocked
assurance snapshot after its canonical pipeline and final PDF complete. This is a
useful systems distinction: derived quality metadata must be reconciled with the
authoritative Stage state before it can serve as a public certificate.

\subsection{Persistent State and Verification-Gated Runtime Self-Evolution}

The long-lived object in \systemname is the campaign state rather than a provider
transcript. Accepted results, failed routes, tools, skills, and current blockers
remain available after session and process restarts. Equation~\ref{eq:process-data}
formalizes the information advantage of this record, while
Equation~\ref{eq:process-compression} explains why the runtime keeps both a complete
event history and a bounded working checkpoint.

This mechanism gives concrete meaning to \emph{verification-gated fixed-model runtime
self-evolution}. The base model remains unchanged, but an admitted mission can
change the starting point, available procedures, verification obligations, or
routing policy of later work. The admission gate binds each candidate to artifacts,
task-native evidence, and an authorized commit; independent Reviewer judgment is one
gate path rather than a requirement for every low-risk update. Rejected branches
also participate when they are retained as verified exclusions: later missions can
avoid repeating the branch and cite its failed evidence when proposing a pivot.
Evolution therefore occurs in the runtime's admitted state and control policy, not
in model weights or an unconstrained summary of prior conversation.

The two longitudinal views expose complementary parts of this claim. The SWE-Bench
Pro trajectory measures how solve Tokens and active time vary as shared state
accumulates, while the mathematical campaign shows direct semantic reuse: later
missions consume earlier definitions, bounds, rejected routes, and bridge lemmas.
Neither trace isolates a single update mechanism. A matched experiment that replays
the same tasks with frozen versus accumulated runtime state would be required to
attribute a causal gain to memory, Skills, verification, or routing separately.

\subsection{Endogenous Harnessing: When a Plan Becomes a Constraint}
\label{sec:endogenous-harness}

A harness is supposed to bound what a system may do. The instructive failure is when
a system manufactures its own harness out of a decision it made earlier, and then
obeys it. We call this \emph{endogenous harnessing}.

The mechanism follows directly from the two levels in Algorithm~1. A Planner proposes
a strategy under the information available at planning time. That proposal is written
into a mission, and from the first round onward the Engineer treats it as the task and
the Reviewer treats it as the standard the work is judged against. What began as one
agent's hypothesis has become an external constraint on every agent that comes after,
including agents that now know more. The system stops behaving like a researcher
continuously reconsidering the best route and starts behaving like a researcher bound
by a note they wrote before the experiment ran.

A recorded verification trace shows the shape. The Planner required a
\code{skip-zero} candidate. The Engineer satisfied that local objective. The Reviewer
then identified a no-gap validator alternative that would have discharged the same
obligation more directly. The Reviewer, however, is scoped to the current round: it
can reject, request revision, or escalate, but it cannot redefine the mission that
made the weaker route mandatory. The better judgement arrived, was recorded, and did
not change the task. The earlier, less informed decision won on authority rather than
on evidence.\footnote{This trace is drawn from internal system-verification work. We
report it as a described sequence rather than a public reproducible artifact.}

Planning is not the problem here, and neither is bounding. The problem is category
confusion between two things a plan can be. A plan can be a \emph{contract}, which is
binding and expensive to change, or a \emph{falsifiable hypothesis}, which is
supposed to be discarded the moment the evidence turns. \systemname currently records
both in the same place, so a technical strategy inherits the immovability that should
belong only to a safety boundary. The correct division is sharper than the one the
runtime implements: frozen authority, prohibitions on granting new trust, and
irreversibility limits are hard harness and should resist revision; the choice of
route, validator, or representation is a technical bet and should be revisable by any
role that produces evidence against it.

The implemented contract already encodes half of this distinction. As
Section~\ref{sec:method} describes, semantic clauses move freely while precise
clauses and the objective require explicit confirmation. What is missing is that a
mission's technical strategy is filed on the immovable side of that line rather than
the negotiable one, and that the Reviewer, which is the role most likely to hold
disconfirming evidence, has no channel to move it.

This reframes where the current bottleneck sits. The failures above are not reasoning
failures. The system found the counterexample, identified the validator alternative,
and understood the cascading invariant. It could not act on any of them in time
because the workflow gave later, better-informed roles no authority to revise the task
boundary. The binding constraint on this class of system is methodological rather than
cognitive, which is also why the remedy is a change in authority routing rather than a
stronger model.

\subsection{From Runtime Learning to Model Training}

The same trajectories can be transformed into supervised examples, preference
pairs, and reinforcement-learning transitions. AgentInstruct, Agent-FLAN, and
Agent Lightning provide complementary mechanisms for this conversion
~\citep{mitra2024agentinstruct,chen2024agentflan,luo2025agentlightning}.
\systemname contributes a structured source stream linking objectives, tool actions,
measurements, revisions, and retained-state updates. A future training stage can use
this stream to internalize recurring planning and verification patterns while the
runtime continues to provide long-horizon coordination.

\subsection{Vertical Research in Mathematics}

The Erd\H{o}s--Gy\'arf\'as campaign exposes the four-role organization at research
depth. The Manager preserves the objective and stage, the Planner converts the
current frontier into bounded theorem tasks, the Engineer produces proofs and
computational checks, and the Reviewer controls admission into durable state.

The campaign retains one falsified route and six accepted theorem-frontier updates.
Later missions improve a bound and close two graph-realizability gaps by starting
from the reviewed ledger rather than restarting from the original conjecture. This
is the same runtime mechanism seen in software repair, applied to a domain where
progress is expressed through definitions, lemmas, counterexamples, and proofs.

\subsection{Verticals Whose Verifier the Runtime Does Not Own}

The chip and materials campaigns (Sections~\ref{sec:silicon-vertical}
and~\ref{sec:materials-vertical}) sharpen the same argument under a stricter
condition: in both, the deciding evaluator is external. A static-timing engine and
an independent chemistry validator will contradict an unsupported claim the moment
they are run, so the gate cannot be satisfied by narrative.

Their outcomes are informative in opposite directions, and both are consistent with
the position of Section~\ref{sec:problem}. ACE-2 closes a scope and then publishes
the perimeter of that scope as part of the result, so the exclusions are an output of
the gate rather than a caveat appended to it. The materials campaign instead reverses
a design assumption: after separating an integrator change from a scoring change, the
runtime found that the published particle-interaction step was discarding the
exploration it paid for, and that keeping $K$ trajectories independent scored higher
at matched compute. The admitted method is therefore smaller than the one it replaced.
Neither outcome required a new component, and neither was reported as a
state-of-the-art claim, because in one case the evidence stops at mapped synthesis and
in the other it stops at the same validator the selection optimizes against.

\subsection{Connecting Theory and Results}

Table~\ref{tab:theory-results} connects the process-to-capability model to the
measurements reported in Sections~\ref{sec:results} and~\ref{sec:vertical-trace}.

\begin{table}[H]
\centering
\scriptsize
\renewcommand{\arraystretch}{1.12}
\begin{tabularx}{\linewidth}{@{}p{2.4cm}p{3.55cm}p{2.55cm}X@{}}
\toprule
\rowcolor{papersand}
\textbf{Theory component} & \textbf{Measurement} & \textbf{Observed result} & \textbf{Interpretation} \\
\midrule
Process-data dominance & Mathematical trajectory & One dead route and six dependent theorem updates & The retained process changes the starting state of later missions. \\
Review gate & 43 Reviewer revision requests & 34 verifier recoveries; 22 strict rescues & Independent review redirects initially unacceptable work into another implementation round. \\
Reuse value $G_L$ & Startup W1--6 versus mature W19--22 & 21\% fewer solve input Tokens; 15\% less active time & Later operation reaches a lower resource regime while shared state expands. \\
Recovery under persistent state & Paper-production trajectory & Seven approaches dropped, two late submission repairs, and 29 session rolls & The campaign changes hypothesis and repairs evidence without losing accepted results or restarting the manuscript. \\
Token conversion & Rounds 12--17 role attribution & 56.0\% reasoning, 55.4\% verification, and 55.6\% action direct shares & The research trace separates frontier work, coordination, and interrupted work. \\
\bottomrule
\end{tabularx}
\caption{Connection between the process-to-capability theory and measured system
behavior. Detailed substitutions appear in Appendix~\ref{app:formula-instantiation}.}
\label{tab:theory-results}
\end{table}
